\chapter{Caractéristiques du flux de transactions}\label{ch:rs:trade-flow-features}

Le 6 mai 2010, le Dow Jones Industrial Average enregistra la plus forte baisse en points de son histoire sur une seule journée, 998.5 points, et, pendant quelques
minutes, mille milliards de dollars de valeur de marché s'évanouirent. Dans l'année qui suivit, trois chercheurs soutinrent qu'une mesure de \omterm{def:rs:trade-flow-features:vpin}{toxicité du flux} d'ordres, le \omterm{def:rs:trade-flow-features:vpin}{VPIN},
avait atteint des niveaux historiquement élevés avant le krach et pouvait aider à signaler la prochaine turbulence. En 2014, une autre étude constata que le \omterm{def:rs:trade-flow-features:vpin}{VPIN} était un mauvais
prédicteur de la volatilité à court terme. Le flux de transactions est la trace laissée par le marché de ceux qui voulaient négocier dans l'urgence, et à quel prix~; en faire des
variables explicatives soulève trois questions auxquelles ce chapitre répond tour à tour~: quel côté a initié chaque transaction, comment les signes des transactions dépendent
les uns des autres, et quand le flux est informé. Ses données sont trois heures de \texttt{firm.tape}, avec un épisode de dix minutes au cours duquel le
prix efficient se déplace huit fois plus vite et où les ordres informés passent de 15~\% à 45~\% du flux~: un épisode toxique dont la vérité est enregistrée par le
simulateur, de sorte que chaque mesure peut être vérifiée par rapport à elle.

\section{Signer les transactions}

\begin{definition}[Signe de la transaction, règle du tick, règle de la cotation, algorithme de Lee--Ready]\label{def:rs:trade-flow-features:sign}
Le \emph{signe de la transaction}\index{signe de la transaction} (trade sign) vaut $+1$ lorsque la transaction a été initiée par un acheteur (un ordre qui a pris le cours vendeur) et $-1$ lorsque
elle a été initiée par un vendeur. La \emph{règle du tick}\index{règle du tick} (tick rule) attribue $+1$ à une transaction dont le prix est supérieur au dernier prix de transaction
différent, et $-1$ s'il est inférieur. La \emph{règle de la cotation}\index{règle de la cotation} (quote rule) lui attribue $+1$ si son prix est supérieur au prix milieu en vigueur, et $-1$
s'il est inférieur. L'\emph{algorithme de Lee--Ready}\index{algorithme de Lee--Ready} applique la règle de la cotation et recourt à la règle du tick
pour les transactions au prix milieu.
\end{definition}

La plupart des données publiques de transactions ne disent pas qui a initié une transaction, et les règles en déduisent l'information. Lee et Ready (1991) ont identifié les deux
problèmes de la \omterm{def:rs:trade-flow-features:sign}{règle de la cotation}~: les cotations peuvent être enregistrées avant les transactions qu'elles ont déclenchées, et les transactions à l'intérieur de la fourchette ne se classent pas
facilement. Le temps est l'ennemi pratique. Sur la bande simulée, où chaque transaction a lieu au meilleur cours, la \omterm{def:rs:trade-flow-features:sign}{règle de la cotation}
avec les cotations en vigueur au moment de la transaction, donne le bon signe pour 100~\% des transactions~; avec des cotations périmées de 0.2 seconde, 99.1~\%~; avec une seconde
de retard, 96.3~\% (\omterm{def:rs:trade-flow-features:sign}{algorithme de Lee--Ready}~: 97.1~\%)~; la \omterm{def:rs:trade-flow-features:sign}{règle du tick}, qui n'a pas besoin de cotations, 95.2~\%
(\cref{fig:rs:trade-flow-features:signing}). Les données réelles font moins bien, car les transactions réelles ont lieu à l'intérieur de la fourchette, dans des ordres cachés et
lors des enchères.

\begin{omfigure}
\begin{tikzpicture}[font=\small,node distance=6mm and 10mm,
  box/.style={draw=omDef,fill=omDef!8,rounded corners=2pt,align=center,inner sep=4pt,minimum height=8mm},
  res/.style={draw=omThm,fill=omThm!10,rounded corners=2pt,align=center,inner sep=4pt,minimum height=8mm},
  arr/.style={-Stealth,thick,draw=black!60}]
\node[box] (q) {cotations à la date de la transaction\\(avec le décalage d'horloge retenu)};
\node[box,below=of q] (c) {prix par rapport au milieu};
\node[res,left=14mm of c] (buy) {achat ($+1$)};
\node[res,right=14mm of c] (sell) {vente ($-1$)};
\node[box,below=of c] (t) {au milieu~: prix par rapport au dernier\\prix de transaction différent};
\node[res,left=14mm of t] (tb) {hausse~: achat};
\node[res,right=14mm of t] (ts) {baisse~: vente};
\draw[arr] (q) -- (c);
\draw[arr] (c) -- node[above,font=\scriptsize]{au-dessus} (buy);
\draw[arr] (c) -- node[above,font=\scriptsize]{en dessous} (sell);
\draw[arr] (c) -- node[right,font=\footnotesize]{égal} (t);
\draw[arr] (t) -- (tb);
\draw[arr] (t) -- (ts);
\end{tikzpicture}

\omcaption{L'\omterm{def:rs:trade-flow-features:sign}{algorithme de Lee--Ready}~: d'abord la \omterm{def:rs:trade-flow-features:sign}{règle de la cotation}, puis la \omterm{def:rs:trade-flow-features:sign}{règle du tick} pour les transactions au milieu. La première case porte l'essentiel de l'erreur~: sur
la bande simulée, un décalage d'horloge d'une seconde coûte à la \omterm{def:rs:trade-flow-features:sign}{règle de la cotation} 3.7 points de précision.}\label{fig:rs:trade-flow-features:leeready}
\end{omfigure}

\begin{definition}[Classification par volume agrégé]\label{def:rs:trade-flow-features:bvc}
La \emph{classification par volume agrégé}\index{classification par volume agrégé} (bulk volume classification, BVC) répartit le volume d'une barre entre achats et ventes sans
signer les transactions une à une~: la part d'achat vaut $\Phi(\Delta P/\sigma_{\Delta P})$, où $\Delta P$ est la variation de prix de la barre et
$\sigma_{\Delta P}$ son écart type sur l'ensemble des barres.
\end{definition}

La BVC a été introduite avec le \omterm{def:rs:trade-flow-features:vpin}{VPIN} pour des données trop grossières pour signer transaction par transaction. Elle a un coût~: sur des barres d'une minute de la bande simulée
elle classe correctement 80.4~\% du volume, contre 95~\% ou plus pour les règles au niveau de la transaction, et Chakrabarty, Pascual et Shkilko
ont constaté, sur un large échantillon d'actions américaines, que la \omterm{def:rs:trade-flow-features:sign}{règle du tick} et l'\omterm{def:rs:trade-flow-features:sign}{algorithme de Lee--Ready} classent nettement mieux que la BVC.

\begin{omfigure}
\begin{tikzpicture}
\begin{axis}[width=11cm,height=5.4cm,xbar,bar width=8pt,xlabel={part classée correctement},tick label style={font=\small},
  label style={font=\small},xmin=0.75,xmax=1.01,ymin=-0.6,ymax=6.6,ytick={0,1,2,3,4,5,6},
  yticklabels={{cotation, à l'heure},{cotation, 0.2 s de retard},{Lee--Ready, 0.2 s de retard},{cotation, 1 s de retard},{Lee--Ready, 1 s de retard},{règle du tick},{BVC, barres d'une minute}},
  y dir=reverse,grid=major,grid style={gray!25},table/col sep=comma]
\addplot[fill=omDef!55,draw=omDef] table[x=accuracy,y=k]{figdata/research/09-trade-flow-features/signing.csv};
\end{axis}
\end{tikzpicture}

\omcaption{Part des exécutions de la bande simulée signées correctement par chaque règle (en volume pour la BVC), par rapport aux
vrais signes du simulateur. Les cotations périmées coûtent à la \omterm{def:rs:trade-flow-features:sign}{règle de la cotation} 3.7 points à une seconde~; la BVC, sur des barres d'une minute, obtient 80~\% du volume correct. Données~:
\texttt{firm.tape}, graine~9.}\label{fig:rs:trade-flow-features:signing}
\end{omfigure}

\section{Autocorrélation du signe des ordres}

\begin{definition}[Autocorrélation du signe des ordres, métaordre]\label{def:rs:trade-flow-features:acf}
L'\emph{autocorrélation du signe des ordres}\index{autocorrélation du signe des ordres} (order-sign autocorrelation) est la fonction d'autocorrélation des signes des
ordres agressifs successifs. Un \emph{métaordre}\index{métaordre} est un ordre de grande taille exécuté sous la forme d'une suite de plus petits ordres enfants répartis dans le temps.
\end{definition}

Les signes des ordres ne sont pas indépendants~: les achats suivent les achats. Lillo et Farmer (2004) ont constaté, à la Bourse de Londres, que
l'autocorrélation des signes des ordres décroît à peu près comme une puissance du retard, d'exposant 0.6, avec un exposant de Hurst de 0.7 (livre~4,
chapitre~17)~: une mémoire longue. L'explication qu'ils ont donnée, avec Mike, par la suite, tient à la fragmentation des \omterm{def:rs:trade-flow-features:acf}{métaordres}.

\begin{proposition}[Les métaordres rendent les signes des ordres à mémoire longue]\label{prop:rs:trade-flow-features:lmf}
Si les \omterm{def:rs:trade-flow-features:acf}{métaordres} sont découpés en ordres enfants de taille constante, si tout le signe du \omterm{def:rs:trade-flow-features:acf}{métaordre} est conservé, et si la distribution des tailles de \omterm{def:rs:trade-flow-features:acf}{métaordres} admet une
queue en loi de puissance $\P(V > v) \propto v^{-\alpha}$ avec $1 < \alpha < 2$, alors l'autocorrélation des signes des ordres enfants décroît asymptotiquement comme
$\tau^{-(\alpha - 1)}$, et l'exposant de Hurst de la série des signes vaut $(3 - \alpha)/2$.
\end{proposition}

\admitted
\noindent Le résultat est dû à Lillo, Mike et Farmer (2005)~; l'exposant de Hurst se déduit de l'exposant de décroissance $\gamma = \alpha - 1$ par
$H = 1 - \gamma/2$.

Les traders de bruit du simulateur tirent les longueurs des \omterm{def:rs:trade-flow-features:acf}{métaordres} d'une loi de Pareto d'exposant de queue 1.5, si bien que la proposition prédit $H = 0.75$. Les
7\,273 ordres agressifs de la bande simulée ont des autocorrélations du signe de 0.20 au retard 1, 0.087 au retard 10 et 0.028 au retard 100
(\cref{fig:rs:trade-flow-features:acf}), et un exposant de Hurst par variance agrégée de 0.715. La mémoire des signes est la raison pour laquelle le volume signé est
un prédicteur des signes suivants, et non seulement un résumé du passé~: un \omterm{def:rs:trade-flow-features:acf}{métaordre} d'achat déjà amorcé continuera probablement d'acheter.

Quelle part de cette mémoire est une prévision des prix~? Le déséquilibre du flux signé des $N$ derniers ordres agressifs (la première \omterm{def:rs:anatomy-of-a-predictor:card}{fiche de prédicteur} ci-dessous) a été
mesuré par rapport à trois cibles sur la bande simulée~: le signe de l'ordre suivant, la variation du prix milieu sur les dix secondes suivantes, et, comme
témoin, la variation sur les dix secondes précédentes.

\begin{center}
\small
\begin{tabular}{@{}rcccc@{}}
\toprule
$N$ derniers ordres & corr.\ signe suivant & signe suivant correct & corr.\ 10 s suivantes & corr.\ 10 s précédentes\\
\midrule
5 & 0.25 & 56.1~\% & 0.16 & 0.22\\
20 & 0.24 & 58.4~\% & 0.14 & 0.31\\
100 & 0.18 & 57.1~\% & 0.12 & 0.17\\
\bottomrule
\end{tabular}
\end{center}

Le flux prévoit bien le signe suivant et moins bien la variation de prix suivante, et pour vingt ordres il est deux fois plus corrélé à la
variation déjà survenue qu'à celle à venir~: une grande partie de ce que le flux sait, le prix le montre déjà. Les signes des ordres sont
prévisibles sans que les prix le soient, ce qui est le propos de Lillo et Farmer~: la liquidité de l'autre côté s'ajuste.

\begin{omfigure}
\begin{tikzpicture}
\begin{loglogaxis}[width=9.5cm,height=5.2cm,xlabel={retard (ordres)},ylabel={autocorrélation du signe des ordres},tick label style={font=\small},
  label style={font=\small},xmin=0.8,xmax=120,ymin=0.01,ymax=0.4,grid=major,grid style={gray!25},legend columns=2,
  legend style={font=\footnotesize,at={(0.5,-0.32)},anchor=north,draw=none,fill=none,
  /tikz/every even column/.append style={column sep=8pt}},table/col sep=comma]
\addplot[omDef,very thick,mark=*] table[x=lag,y=acf]{figdata/research/09-trade-flow-features/signacf.csv};
\addplot[omThm,thick,dashed,domain=1:100,samples=2] {0.2*x^(-0.5)};
\legend{tape simulé,$0.2\,\tau^{-0.5}$ (pente de la proposition)}
\end{loglogaxis}
\end{tikzpicture}

\omcaption{Autocorrélation des signes des ordres agressifs successifs sur le tape simulé, en axes logarithmiques, avec une loi de puissance de
exposant $0.5 = \alpha - 1$ pour la queue des \omterm{def:rs:trade-flow-features:acf}{métaordres} du simulateur $\alpha = 1.5$. Données~: \texttt{firm.tape}, graine 9.}\label{fig:rs:trade-flow-features:acf}
\end{omfigure}

\section{Toxicité du flux}

\begin{definition}[Toxicité du flux, VPIN]\label{def:rs:trade-flow-features:vpin}
La \emph{toxicité du flux}\index{toxicité du flux} est la sélection adverse que le flux d'ordres impose aux fournisseurs de liquidité qui négocient contre lui~:
ses agresseurs sont informés, et le côté passif perd en moyenne. Le \emph{VPIN}\index{VPIN} (probabilité de négociation informée synchronisée sur le
volume) découpe la suite des transactions en compartiments de même volume $V$ et, à la fin de chaque compartiment, moyenne $|V^B_\tau - V^S_\tau|/V$ sur
les $n$ derniers compartiments, $V^B_\tau$ et $V^S_\tau$ désignant le volume acheteur et le volume vendeur du compartiment.
\end{definition}

Le \omterm{def:rs:trade-flow-features:vpin}{VPIN} mesure un flux unilatéral. Ses auteurs définissent la toxicité comme la sélection adverse des teneurs de marché (Easley, López de Prado et O'Hara,
2012) et soutiennent que c'est par le volume unilatéral qu'elle se manifeste. La toxicité elle-même, en revanche, se mesure directement a posteriori~: le mark-out
du livre~2, chapitre~15, des transactions contre leur côté passif. L'épisode simulé sépare nettement les deux. Pendant celui-ci, le
mark-out des agresseurs dix secondes après la transaction passe de $-0.36$ ticks (ils paient la demi-fourchette) à $+0.96$ (ils gagnent)~: le flux est toxique par
définition. Le \omterm{def:rs:trade-flow-features:vpin}{VPIN}, sur des compartiments de 2\,649 actions et une fenêtre de vingt, passe de 0.31 à 0.35, soit une hausse de 13~\%~; son maximum sur les trois heures
se situe hors de l'épisode, 20 minutes après sa fin (la \cref{fig:rs:trade-flow-features:vpin}). Les traders informés du simulateur achètent lorsque
le prix efficient est au-dessus du prix milieu et vendent lorsqu'il est en dessous, si bien que leur flux n'est pas unilatéral sur un compartiment~; ce sont les \omterm{def:rs:trade-flow-features:acf}{métaordres} des traders de bruit
qui le sont. Le volume unilatéral et le volume informé sont deux choses différentes.

\begin{omfigure}
\begin{tikzpicture}
\begin{axis}[name=top,scale only axis,width=9.6cm,height=3.2cm,ylabel={VPIN},tick label style={font=\small},label style={font=\small},xmin=0,xmax=180,
  ymin=0.15,ymax=0.6,grid=major,grid style={gray!25},xticklabels={},table/col sep=comma]
\fill[omThm!12] (axis cs:90,0.15) rectangle (axis cs:100,0.6);
\addplot[omDef,thick] table[x=minute,y=vpin]{figdata/research/09-trade-flow-features/vpin.csv};
\end{axis}
\begin{axis}[at={(top.south west)},anchor=north west,yshift=-3mm,scale only axis,width=9.6cm,height=3.2cm,xlabel={minutes},
  ylabel={mark-out (en ticks)},tick label style={font=\small},label style={font=\small},xmin=0,xmax=180,ymin=-1.5,ymax=3.2,grid=major,
  grid style={gray!25},table/col sep=comma]
\fill[omThm!12] (axis cs:90,-1.5) rectangle (axis cs:100,3.2);
\addplot[omThm,thick] table[x=minute,y=markout]{figdata/research/09-trade-flow-features/markout.csv};
\end{axis}
\end{tikzpicture}

\omcaption{L'épisode toxique (grisé, minutes 90 à 100) vu par le \omterm{def:rs:trade-flow-features:vpin}{VPIN} (haut~: 400 compartiments de 2\,649 actions, fenêtre de 20) et par le
mark-out des agresseurs dix secondes après chaque transaction (bas~: moyenne mobile sur une minute, tracée lorsqu'elle est connue). Le mark-out devient fortement
positif~; le \omterm{def:rs:trade-flow-features:vpin}{VPIN} bondit au début de l'épisode, mais sa moyenne sur l'épisode ne progresse que de 13~\% et son maximum se situe hors de celui-ci. Données~: \texttt{firm.tape}, graine 9.}\label{fig:rs:trade-flow-features:vpin}
\end{omfigure}

Le débat sur le flash crash avait la même forme. Easley, López de Prado et O'Hara (2011) ont relevé des valeurs historiquement élevées du \omterm{def:rs:trade-flow-features:vpin}{VPIN} avant le
crash~; Andersen et Bondarenko (2014) ont jugé le \omterm{def:rs:trade-flow-features:vpin}{VPIN} mauvais prédicteur de la volatilité à court terme. Un simulateur ne peut pas trancher un débat sur les
marchés réels, mais il peut montrer qu'une mesure échoue sur la toxicité même qu'elle était censée détecter, ce qui suffit à justifier de la valider sur des données où la
vérité est connue avant de lui faire confiance là où elle ne l'est pas. La mesure par mark-out a son propre coût~: elle exige le futur (dix secondes ici), et elle
détecte donc tardivement~; dans l'épisode, elle franchit pour la première fois son 95e centile d'avant l'épisode 40 secondes après le début, le \omterm{def:rs:trade-flow-features:vpin}{VPIN} au bout de 16 secondes, mais le \omterm{def:rs:trade-flow-features:vpin}{VPIN}
passe aussi 6.0~\% du reste de la séance au-dessus de son seuil.

\section{Intensités d'excitation ajustées}

Les ordres agressifs se regroupent dans le temps, et un processus de Hawkes (livre~4, chapitre~7) en est la description standard~: chaque ordre fait monter l'intensité
de l'ordre suivant. Son ratio de branchement est la part des ordres déclenchés par les précédents. Sur une version plate du tape simulé (sans processus d'activité
ni nouvelles, ni traders informés), l'estimation par maximum de vraisemblance retrouve 0.403, contre les 0.4 placés par le simulateur. Sur le
tape du chapitre, l'ajustement donne 0.65. Rien dans le mécanisme de génération des ordres n'a changé~; ce qui a changé, c'est que l'activité monte et descend pour d'autres
raisons, et un processus de Hawkes attribue chaque grappe à l'auto-excitation. Un ratio de branchement ajusté est une borne supérieure de l'endogénéité
sauf si l'on laisse la base varier. Comme attributs, les intensités ajustées des achats et des ventes, et leur différence, résument le degré d'excitation de chaque
côté du flux.

\section{Grandes transactions et métaordres}

\begin{definition}[Lambda de Kyle]\label{def:rs:trade-flow-features:lambda}
Le \emph{lambda de Kyle}\index{lambda de Kyle} est l'impact sur le prix par unité de flux d'ordres signé~: la pente des variations de prix sur le volume net signé
sur les mêmes intervalles, au sens du paramètre du modèle de Kyle (1985) d'un trader informé et des teneurs de marché.
\end{definition}

Sur le tape simulé, sur des intervalles de dix secondes, le lambda vaut 0.10 tick pour 100 actions d'achat net, avec un $R^2$ de 14~\%~: très en deçà des
53~\% du \omterm{def:rs:order-book-features:ofi}{déséquilibre du flux d'ordres} (chapitre~8), car le volume signé des transactions ne tient pas compte des annulations et des ajouts qui font bouger les cotations, comme l'ont souligné Cont,
Kukanov et Stoikov pour des données réelles. Le lambda est une caractéristique de liquidité (une action dont le prix bouge davantage par action échangée est moins liquide)
et une donnée d'entrée pour le coût (chapitre~27). La détection des \omterm{def:rs:trade-flow-features:acf}{métaordres} part des mêmes flux~: une suite d'ordres de même signe et de taille voisine à intervalles
réguliers, une part inhabituelle d'un côté sur quelques minutes, un déséquilibre de type \omterm{def:rs:trade-flow-features:vpin}{VPIN} sur une horloge de volume. Le résultat de mémoire du chapitre dit pourquoi cela
fonctionne~: les enfants d'un même parent partagent son signe.

\section{Fiches de prédicteur}

\begin{predictorcard}{Déséquilibre du flux signé}\label{pred:rs:trade-flow-features:signed}
\sfield{Définition} $\sum b_iq_i/\sum q_i$ sur les ordres agressifs des $\Delta$ dernières secondes (ou des $N$ derniers ordres), $b_i$ désignant le signe.
\sfield{Entrées et horodatage} Transactions avec heures de réception~; signes issus de l'indicateur d'agresseur du flux, ou de l'\omterm{def:rs:trade-flow-features:sign}{algorithme de Lee--Ready} avec les cotations en vigueur à l'instant
de réception de la transaction.
\sfield{Justification} Les signes des ordres ont une mémoire longue (la \cref{prop:rs:trade-flow-features:lmf})~; les \omterm{def:rs:trade-flow-features:acf}{métaordres} se poursuivent.
\sfield{Horizon et demi-vie} Les prochains ordres~; \omterm{def:rs:trade-flow-features:acf}{autocorrélation du signe des ordres} de 0.20 à un ordre, 0.087 à dix sur le tape simulé~; sur vingt ordres, il prédit le signe suivant juste dans 58.4~\% des cas.
\sfield{Normalisation} Par volume total~; d'un instrument à l'autre, par taille moyenne des transactions.
\sfield{Modes de défaillance} Mauvaise classification avec des cotations périmées~; le prix bouge avant que le flux soit observé (l'impact est déjà dans le prix).
\sfield{Sources} Lillo et Farmer (2004)~; Lillo, Mike et Farmer (2005)~; \texttt{rs\_tradeflow.memory}.
\end{predictorcard}

\begin{predictorcard}{VPIN}\label{pred:rs:trade-flow-features:vpin}
\sfield{Définition} Déséquilibre absolu moyen entre achats--ventes sur les $n$ derniers compartiments de volume, divisé par la taille du compartiment.
\sfield{Entrées et horodatage} Transactions signées (ou BVC sur les barres), découpées en compartiments de même volume.
\sfield{Justification} Le flux toxique est unilatéral (Easley, López de Prado et O'Hara, 2012).
\sfield{Horizon et demi-vie} De quelques minutes à quelques heures. Sur l'épisode simulé, il progresse de 13~\% tandis que le mark-out des agresseurs passe de $-0.36$
à $+0.96$ ticks.
\sfield{Normalisation} Borné dans $[0, 1]$~; son niveau dépend de la taille des compartiments et de la fenêtre.
\sfield{Modes de défaillance} Flux informé qui n'est pas unilatéral~; flux unilatéral qui n'est pas informé (\omterm{def:rs:trade-flow-features:acf}{métaordres}, rééquilibrages d'indices)~; mauvaise
prévision de la volatilité (Andersen et Bondarenko, 2014).
\sfield{Sources} Comme cité~; \texttt{rs\_tradeflow.toxicity\_race}.
\end{predictorcard}

\section{Tutoriel~: signer et pondérer le flux}\label{tut:rs:trade-flow-features:flow}

\begin{tutorial}
\textbf{Objectif.} Signer les transactions de la bande simulée selon chaque règle, mesurer la mémoire du signe des ordres et ajuster un processus de Hawkes, et confronter le \omterm{def:rs:trade-flow-features:vpin}{VPIN}
au mark-out tout au long de l'épisode toxique. \textbf{État final~:}
les \cref{fig:rs:trade-flow-features:signing,fig:rs:trade-flow-features:acf,fig:rs:trade-flow-features:vpin}~; les ratios de branchement de Hawkes valent 0.65
et 0.403.
\begin{tutsteps}
\item \textbf{Les règles de signature.} La \omterm{def:rs:trade-flow-features:sign}{règle du tick} propage le dernier changement de prix~; l'\omterm{def:rs:trade-flow-features:sign}{algorithme de Lee--Ready} utilise la \omterm{def:rs:trade-flow-features:sign}{règle de la cotation} et la \omterm{def:rs:trade-flow-features:sign}{règle du tick} pour
les transactions au prix milieu.
\omcode{code/firm/tradeflow/firm_tradeflow.py}{27}{49}{Règle du tick, règle de la cotation et Lee--Ready.}\label{lst:rs:trade-flow-features:rules}
\item \textbf{\omterm{def:rs:trade-flow-features:vpin}{VPIN}.} Les transactions sont versées dans des compartiments de même volume, coupées là où une transaction est à cheval sur deux compartiments.
\omcode{code/firm/tradeflow/firm_tradeflow.py}{79}{101}{VPIN sur compartiments de volume.}\label{lst:rs:trade-flow-features:vpincode}
\item \textbf{Lancer} \texttt{signing()}, \texttt{memory()}, \texttt{signed\_flow()}, \texttt{hawkes\_fits()}, \texttt{toxicity\_race()}, \texttt{kyle()} et
\texttt{fig\_tradeflow.py}.
\end{tutsteps}
\textbf{Pour aller plus loin.} Rendre les traders informés unilatéraux (n'acheter que tant que dure l'épisode) et observer la réaction du \omterm{def:rs:trade-flow-features:vpin}{VPIN}~; ajuster un processus de
Hawkes dont la base suit le niveau d'activité et observer la redescente du ratio de branchement.
\end{tutorial}

\section{Construction~: la boîte à outils du flux de transactions}\label{bld:rs:trade-flow-features:tradeflow}

\begin{build}
\bfield{Rôle} Mesures de signature, de mémoire, de toxicité et d'impact sur les données de transactions de la firme miniature, pour les caractéristiques (les livres~8, 11 et~12) et pour
la recherche en exécution (livre~10).
\bfield{Interface} \texttt{tick\_rule}, \texttt{quote\_rule}, \texttt{lee\_ready}, \texttt{bvc}, \texttt{sign\_acf}, \texttt{hurst\_aggvar},
\texttt{vpin(buy\_volume, volume, bucket, n)}, \texttt{kyle\_lambda}, \texttt{markout\_toxicity}, \texttt{aggregate\_orders}.
\bfield{Règles} Les cotations transmises sont celles qui prévalent au \omterm{def:rs:point-in-time-data-and-the-biases:bitemporal}{temps de connaissance} de chaque transaction~; une transaction à cheval sur un compartiment de volume est répartie
au prorata~; les exécutions d'un même ordre agressif sont agrégées avant les statistiques de signe.
\bfield{Tests d'acceptation} \texttt{code/firm/tradeflow/tests/}~: les trois règles sur une séquence construite à la main et la symétrie de la BVC~; le \omterm{def:rs:trade-flow-features:vpin}{VPIN}
d'une séquence équilibrée et d'une séquence unilatérale~; $H = 0.5$ pour des données indépendantes et une autocorrélation du signe qui s'annule au-delà de la longueur d'une série~; un
\omterm{def:rs:trade-flow-features:lambda}{lambda de Kyle} connu retrouvé~; agrégation des ordres.
\bfield{Pour aller plus loin} Ajustements de Hawkes avec une base variable dans le temps~; détection des \omterm{def:rs:trade-flow-features:acf}{métaordres} par points de rupture dans le flux signé~; mark-outs par
type de contrepartie.
\end{build}

\begin{omsources}
\item C.~M.~C.~Lee and M.~J.~Ready, ``Inferring trade direction from intraday data'', \emph{Journal of Finance} 46(2), 1991.
\item B.~Chakrabarty, R.~Pascual and A.~Shkilko, ``Evaluating trade classification algorithms: bulk volume classification versus the tick rule
and the Lee--Ready algorithm'', \emph{Journal of Financial Markets} 25, 2015.
\item F.~Lillo and J.~D.~Farmer, ``The long memory of the efficient market'', \emph{Studies in Nonlinear Dynamics and Econometrics} 8(3), 2004;
F.~Lillo, S.~Mike and J.~D.~Farmer, ``Theory for long memory in supply and demand'', \emph{Physical Review E} 71, 2005.
\item D.~Easley, M.~López de Prado and M.~O'Hara, ``The microstructure of the `flash crash''', \emph{Journal of Portfolio Management}
37(2), 2011; ``Flow toxicity and liquidity in a high-frequency world'', \emph{Review of Financial Studies} 25(5), 2012.
\item T.~G.~Andersen and O.~Bondarenko, ``VPIN and the flash crash'', \emph{Journal of Financial Markets} 17, 2014.
\item A.~S.~Kyle, ``Continuous auctions and insider trading'', \emph{Econometrica} 53(6), 1985.
\end{omsources}

\section{Exercices}

\begin{exercise}[$\star$]\label{exo:rs:trade-flow-features:1}
Transactions à 20.00, 20.01, 20.01, 20.00, 19.99 avec les cotations 20.00--20.01 pendant toute la période. Signer chacune par la \omterm{def:rs:trade-flow-features:sign}{règle du tick} et par Lee--Ready.
\end{exercise}

\begin{exercise}[$\star$]\label{exo:rs:trade-flow-features:2}
La variation de prix d'une barre d'une minute vaut $+0.8$ écart-type. Quelle part de son volume la BVC classe-t-elle en achats~?
\end{exercise}

\begin{exercise}[$\star$]\label{exo:rs:trade-flow-features:3}
\omterm{def:rs:trade-flow-features:acf}{Métaordres} dont l'exposant de queue des tailles vaut 1.4~: quels exposants de décroissance et de Hurst la \cref{prop:rs:trade-flow-features:lmf} prédit-elle~?
\end{exercise}

\begin{exercise}[$\star\star$]\label{exo:rs:trade-flow-features:4}
Cinq compartiments de volume de 1\,000 actions ont des volumes d'achat de 700, 400, 900, 500 et 200. Quel est le \omterm{def:rs:trade-flow-features:vpin}{VPIN} sur les cinq derniers compartiments~?
\end{exercise}

\begin{exercise}[$\star\star$]\label{exo:rs:trade-flow-features:5}
Le \omterm{def:rs:trade-flow-features:lambda}{lambda de Kyle} vaut 0.10 tick pour 100 actions. Quelle variation de prix un achat net de 3\,000 actions sur dix secondes implique-t-il, et pourquoi
l'estimation est-elle plus bruitée que celle fondée sur le \omterm{def:rs:order-book-features:ofi}{déséquilibre du flux d'ordres}~?
\end{exercise}

\begin{exercise}[$\star\star$]\label{exo:rs:trade-flow-features:6}
Pourquoi un ajustement de Hawkes sur des données dont l'activité de base varie surestime-t-il le ratio de branchement~? Donner l'intuition avec deux régimes
d'intensité constante différente.
\end{exercise}

\begin{exercise}[$\star\star\star$]\label{exo:rs:trade-flow-features:7}
\emph{Programmation.} Avec la bande du chapitre, calculer la part des transactions signées correctement par la \omterm{def:rs:trade-flow-features:sign}{règle de la cotation} avec des cotations en retard de
0.05, 0.5 et 2 secondes. À quelle vitesse la précision chute-t-elle avec l'ancienneté des cotations~?
\end{exercise}

\begin{exercise}[$\star\star\star$]\label{exo:rs:trade-flow-features:8}
\emph{Trouver la faille.} «~Le \omterm{def:rs:trade-flow-features:vpin}{VPIN} a atteint son niveau le plus haut de l'année une heure avant le krach, donc le \omterm{def:rs:trade-flow-features:vpin}{VPIN} prédit les krachs.~»
\end{exercise}

\section{Problème~: le VPIN a-t-il vu venir le coup~?}

\begin{problem}[{Problème du week-end --- une mesure de toxicité face à la vérité}]\label{pb:rs:trade-flow-features:1}
La bande du chapitre~: trois heures, dont un épisode de dix minutes à partir de la minute 90 pendant lequel le prix efficace évolue huit fois plus vite et où
la négociation de bruit ne progresse pas, graine 9.

\medskip\noindent\textbf{Partie I --- La vérité.}
\begin{enumerate}
\item Quelle part du flux est informé avant et pendant l'épisode~?
\item Quel est le mark-out à dix secondes des agresseurs, en ticks, avant et pendant l'épisode~?
\item Pourquoi le mark-out est-il négatif avant l'épisode~?
\item Le flux de l'épisode est-il toxique selon la définition du chapitre~?
\end{enumerate}

\medskip\noindent\textbf{Partie II --- \omterm{def:rs:trade-flow-features:vpin}{VPIN}.}
\begin{enumerate}[resume]
\item Quelles sont la taille des compartiments et la fenêtre retenues par le chapitre, et combien la session compte-t-elle de compartiments~?
\item Quelles sont les moyennes du \omterm{def:rs:trade-flow-features:vpin}{VPIN} avant et pendant l'épisode~?
\item Quand le \omterm{def:rs:trade-flow-features:vpin}{VPIN} atteint-il son maximum~?
\item Quand franchit-il pour la première fois, après le début de l'épisode, son 95e centile d'avant l'épisode, et quelle part du reste de la session
se situe au-dessus de ce seuil~?
\item Pourquoi le \omterm{def:rs:trade-flow-features:vpin}{VPIN} bouge-t-il à peine~?
\end{enumerate}

\medskip\noindent\textbf{Partie III --- La mesure du mark-out.}
\begin{enumerate}[resume]
\item Quelles sont ses moyennes avant et pendant l'épisode, et quand atteint-il son pic~?
\item Quand franchit-il pour la première fois son seuil, et pourquoi est-ce plus tard que le \omterm{def:rs:trade-flow-features:vpin}{VPIN}~?
\item Quelle part du reste de la session se situe au-dessus de son seuil~?
\item Quelle mesure retenir pour protéger un teneur de marché, et laquelle pour étudier la toxicité après coup~?
\end{enumerate}

\medskip\noindent\textbf{Partie IV --- Au-delà.}
\begin{enumerate}[resume]
\item Qu'est-ce qui rendrait le flux informé unilatéral, et rendrait le \omterm{def:rs:trade-flow-features:vpin}{VPIN} utile~?
\item Que donne l'ajustement de Hawkes sur ce tape et sur le tape plat~?
\item Quels sont le \omterm{def:rs:trade-flow-features:lambda}{lambda de Kyle} et son $R^2$ à dix secondes~?
\item Laquelle des caractéristiques du chapitre pourrait détecter l'épisode le plus vite, et à quel prix en fausses alarmes~?
\item Énoncer le \emph{résultat nommé}~: la hausse du \omterm{def:rs:trade-flow-features:vpin}{VPIN} pendant un épisode où la part informée triple et où le mark-out des agresseurs passe
de $-0.36$ à $+0.96$ tick.
\item Que dit le résultat sur la validation d'une mesure de toxicité~?
\item En une phrase~: qu'est-ce que le flux toxique~?
\end{enumerate}
\end{problem}

\section{Questions d'entretien}

\begin{interviewq}[$\star$ \iqroles{chercheur, trader}]\label{iq:rs:trade-flow-features:1}
Comment décideriez-vous si une transaction était un achat ou une vente lorsque les données ne le disent pas~?
\end{interviewq}

\begin{interviewq}[$\star\star$ \iqroles{chercheur}]\label{iq:rs:trade-flow-features:2}
Pourquoi les signes des ordres au marché sont-ils autocorrélés, et pourquoi cela ne rend-il pas les prix prévisibles de la même façon~?
\end{interviewq}

\begin{interviewq}[$\star\star$ \iqroles{trader, chercheur}]\label{iq:rs:trade-flow-features:3}
Qu'est-ce que la \omterm{def:rs:trade-flow-features:vpin}{toxicité du flux}, et comment la mesureriez-vous pour votre propre carnet de teneur de marché~?
\end{interviewq}

\begin{interviewq}[$\star\star$ \iqroles{chercheur, ingénieur ML}]\label{iq:rs:trade-flow-features:4}
Vous ajustez un processus de Hawkes sur les instants de transaction et obtenez un ratio de branchement de 0.9. Qu'est-ce qui d'autre pourrait l'expliquer~?
\end{interviewq}

\begin{interviewq}[$\star\star$ \iqroles{chercheur}]\label{iq:rs:trade-flow-features:5}
Expliquez le \omterm{def:rs:trade-flow-features:vpin}{VPIN} et sa principale critique.
\end{interviewq}

\begin{interviewq}[$\star\star\star$ \iqroles{chercheur}]\label{iq:rs:trade-flow-features:6}
Établissez la relation entre l'exposant de queue des tailles de \omterm{def:rs:trade-flow-features:acf}{métaordres} et l'exposant de Hurst des signes des ordres, puis testez-la sur une série.
\end{interviewq}
